% A \chardef AND ITS KIN MAKE THEIR TARGET \relax BEFORE SCANNING.
%
% The name being defined stands as \relax for the whole of the value scan,
% whatever it meant before and whether or not it meant anything. Two
% things turn on it:
%
%   the value may ask what the name means, and gets \relax;
%   a number that looks ahead past its last digit and meets the name
%   stops there, rather than expanding it -- which for a name that means
%   nothing yet would report "Undefined control sequence" and swallow it.
%
% trip line 160 writes
%
%   \dimendef\varunit=222\varunit=+1,001\ifdim.5\mag>0cc0\fi1pt
%
% where the 222 looks ahead onto \varunit itself. The standing-in \relax
% ends the number, is put back, and is then read again with the meaning
% the \dimendef has by that time given it.
%
% \global reaches the standing-in \relax as it reaches the real meaning.
\catcode`\{=1 \catcode`\}=2
\tracingonline=1
\message{[A a name that meant nothing]}
\chardef\ca=\ifx\ca\relax 5\else 6\fi
\showthe\ca
\message{[B a name that meant something]}
\def\cb{X}
\chardef\cb=\ifx\cb\relax 7\else 8\fi
\showthe\cb
\message{[C a number that looks ahead onto its own target]}
\dimendef\varunit=222\varunit=3pt
\showthe\dimen222 \showthe\varunit
\message{[D the same for countdef and mathchardef]}
\countdef\cn=44\cn=9 \showthe\count44
\mathchardef\cm=\ifx\cm\relax "161\else "162\fi
\showthe\cm
\message{[E globally]}
{\global\chardef\cc=\ifx\cc\relax 1\else 2\fi }
\showthe\cc
\message{[done]}
\end
